% Chapter 0: 引言 - 文献概述
% Tang-Rong 流形统计学习系列

\chapter*{引言：文献概述}
\addcontentsline{toc}{chapter}{引言：文献概述}
\label{chap:Introduction}

\section{总述}

本笔记学习的主题是\textbf{流形统计学习}（Statistical Learning on Manifolds），
由 Rong Tang 与 Yun Yang 等人在 2021--2025 年间系统发展的一系列研究。
该系列涵盖 9 篇论文，发表于 COLT、Annals of Statistics、JRSS B、ICML、AISTATS 等
顶级会议与期刊，并延伸至分布回归与条件扩散模型的 arXiv 工作。

整个系列围绕三个核心问题展开：
\begin{enumerate}
    \item \textbf{统计极限}：当数据支撑于未知低维流形时，分布估计与回归的极小极大速率是什么？
    \item \textbf{自适应机制}：深度生成模型（VAE、扩散模型）能否隐式自适应流形结构？
    \item \textbf{贝叶斯推断}：如何在流形约束下进行稳健的贝叶斯推断和高效 MCMC？
\end{enumerate}

三条研究主线相互呼应，形成从理论奠基到方法验证、再到算法优化的完整体系。

\section{文献摘要总结}\label{sec:LiteratureSummary}

\subsection{Excess Risk Bounds for Distribution Estimation with VAE}
\begin{enumerate}
    \item \textbf{作者:} Rong Tang, Yun Yang
    \item \textbf{会议:} COLT 2021
    \item \textbf{年份:} 2021
    \item \textbf{篇幅:} 58 页
    \item \textbf{核心贡献:}
        \begin{itemize}
            \item 为 VAE 型生成模型提供了首个关于 Excess Risk 的非渐近上界
            \item 揭示了 VAE 隐式自适应低维流形结构的机制——
                  潜在空间维数与流形本征维数的关系是关键
            \item 给出了 VAE 在分布估计任务中的理论保证
        \end{itemize}
    \item \textbf{摘要:} 本文为变分自编码器（VAE）在分布估计任务中的 Excess Risk 建立了非渐近上界。结果表明，当数据支撑于低维流形时，VAE 的估计误差仅依赖流形的本征维度而非环境维度，揭示了深度生成模型隐式学习流形结构的理论机制。
    \item \textbf{关键词:} VAE； Excess Risk； 低维流形； 分布估计； 变分推断
\end{enumerate}

\subsection{Minimax Rate of Distribution Estimation on Unknown Submanifold}
\begin{enumerate}
    \item \textbf{作者:} Rong Tang, Yun Yang
    \item \textbf{期刊:} Annals of Statistics, 2022, Vol.50, No.4, 2197--2216
    \item \textbf{年份:} 2022
    \item \textbf{篇幅:} 102 页
    \item \textbf{核心贡献:}
        \begin{itemize}
            \item 首次建立了未知子流形上分布估计在对抗损失下的完整极小极大理论
            \item 识别了三类收敛机制：参数机制、密度估计机制、流形估计机制
            \item 关键发现：环境维度 $D$ 完全不出现在速率指数中
            \item 引入单位分解（Partition of Unity）技术，构造覆盖无全局参数化流形的估计器
        \end{itemize}
    \item \textbf{摘要:} 本文研究了数据支撑于未知低维子流形时的分布估计问题。在对抗损失（adversarial loss）框架下，建立了完整的极小极大理论，精确刻画了分布估计的收敛速率。核心发现是三类不同收敛机制的识别与竞争，揭示了流形几何参数如何共同决定统计极限。
    \item \textbf{关键词:} 极小极大速率； 未知子流形； 对抗损失； 单位分解； 热核平滑
\end{enumerate}

\subsection{Estimating Distributions with Low-dimensional Structures Using Mixtures of Generative Models}
\begin{enumerate}
    \item \textbf{作者:} Rong Tang, Yun Yang
    \item \textbf{会议:} ICML 2023
    \item \textbf{年份:} 2023
    \item \textbf{篇幅:} 待补充
    \item \textbf{核心贡献:}
        \begin{itemize}
            \item 提出混合潜在分布匹配自编码器（MLDMAE）方法
            \item 利用 $K$ 个局部编码器-解码器对分别学习流形局部区域的参数化
            \item 通过单位分解函数 $\rho_k$ 实现光滑拼接
            \item 在 1-Wasserstein 距离下达极小极大最优（至多对数因子）
        \end{itemize}
    \item \textbf{摘要:} 本文提出 MLDMAE 方法，将 \citep{tang2022minimax} 的理论极小极大框架转化为可计算的深度学习方法。多编码器/解码器结构在流形分布估计中是必要的——单编码器方法无法处理无全局参数化的流形。
    \item \textbf{关键词:} 混合生成模型； 单位分解； Wasserstein 距离； 流形分布估计； 自编码器
\end{enumerate}

\subsection{Robust Bayesian Inference on Riemannian Submanifold}
\begin{enumerate}
    \item \textbf{作者:} Rong Tang, Anirban Bhattacharya, Debdeep Pati, Yun Yang
    \item \textbf{期刊:} Journal of the Royal Statistical Society: Series B, 2023, Vol.85, No.4, 1079--1103
    \item \textbf{年份:} 2023
    \item \textbf{篇幅:} 103 页
    \item \textbf{核心贡献:}
        \begin{itemize}
            \item 提出 Bayesian RPETEL 框架——伪似然 + 经验风险惩罚，无需正确指定似然函数
            \item 建立黎曼流形上的 Bernstein-von Mises（\BvM）定理：后验在切空间投影上渐近正态
            \item 提出 RRWM（Riemannian Random-Walk Metropolis）算法，混合时间仅依赖 $d$ 而非 $D$
            \item 证明三重稳健性：似然稳健性、风险稳健性、先验稳健性
        \end{itemize}
    \item \textbf{摘要:} 本文研究了参数取值于黎曼子流形时的贝叶斯推断问题。即使后验分布在 $\mathbb{R}^D$ 中是奇异的（集中在 $d$ 维子流形上），其在切空间上的投影仍渐近正态。Bayesian RPETEL 框架在模型误设定下展现出三重稳健性，RRWM 算法在高维流形问题中具有显著计算优势。
    \item \textbf{关键词:} 黎曼流形； 贝叶斯推断； Bernstein-von Mises 定理； 稳健性； MCMC
\end{enumerate}

\subsection{Minimax Nonparametric Two-Sample Test under Adversarial Losses}
\begin{enumerate}
    \item \textbf{作者:} Rong Tang, Yun Yang
    \item \textbf{会议:} AISTATS 2023
    \item \textbf{年份:} 2023
    \item \textbf{篇幅:} 34 页
    \item \textbf{核心贡献:}
        \begin{itemize}
            \item 建立了统一的两样本检验框架，基于负 Besov 范数的样本版本
            \item 建立了负 Besov 范数与对抗损失之间的等价关系
            \item 框架同时适用于 $\ell_1$、$\ell_2$ 距离和 Wasserstein 距离等多种度量
            \item 计算复杂度 $O(n\log n + n^{2d/(4\alpha+d)})$，远优于 MMD 的 $O(n^2)$
        \end{itemize}
    \item \textbf{摘要:} 本文研究流形支撑数据的两样本检验问题。通过构造基于截断小波展开的负 Besov 范数，建立了对抗损失下的极小极大最优检验理论。结果表明，热核正则化技术将流形上的奇异分布检验问题转化为标准的密度检验问题。
    \item \textbf{关键词:} 两样本检验； 对抗损失； Besov 范数； 流形数据； 热核平滑
\end{enumerate}

\subsection{Adaptivity of Diffusion Models to Manifold Structures}
\begin{enumerate}
    \item \textbf{作者:} Rong Tang, Yun Yang
    \item \textbf{会议:} AISTATS 2024
    \item \textbf{年份:} 2024
    \item \textbf{篇幅:} 45 页
    \item \textbf{核心贡献:}
        \begin{itemize}
            \item 证明两类扩散模型（Langevin 扩散和正向--反向扩散）的分布估计收敛速率仅依赖本征维度 $d$ 而非环境维度 $D$
            \item 通过 Girsanov 定理建立分数估计误差与分布估计误差之间的联系
            \item 关键改进：只需分数函数的四阶矩误差即可控制分布估计误差
            \item 在 1-Wasserstein 距离下达极小极大最优速率
        \end{itemize}
    \item \textbf{摘要:} 本文建立了扩散模型流形自适应理论的完整框架。核心发现是：正向扩散天然地通过高斯噪声填充了流形的管状邻域，从而回避了奇异分布（流形支撑分布）的困难。条件扩散模型和边际扩散模型均达到极小极大最优收敛速率。
    \item \textbf{关键词:} 扩散模型； 流形自适应； 分数匹配； Girsanov 定理； Wasserstein 距离
\end{enumerate}

\subsection{Metropolis-adjusted Langevin Algorithm: Optimal Dimension Dependency and Beyond}
\begin{enumerate}
    \item \textbf{作者:} Rong Tang, Yun Yang
    \item \textbf{会议:} ICML 2024
    \item \textbf{年份:} 2024
    \item \textbf{篇幅:} 66 页
    \item \textbf{核心贡献:}
        \begin{itemize}
            \item 引入 $s$-conductance profile 新技术，证明 MALA 达到最优 $O(d^{1/3})$ 维度依赖
            \item 无需对目标分布施加光滑或强对数凹性假设
            \item 将 warm-start 参数依赖从 $\log M_0$ 改进到 $\log\log M_0$
            \item 给出了匹配的下界，证明最优性
        \end{itemize}
    \item \textbf{摘要:} 本文研究了 Metropolis-adjusted Langevin 算法（MALA）的混合时间分析。通过引入 $s$-conductance profile 技术，在无需强假设的条件下证明了 MALA 的最优维度依赖。这一结果与黎曼流形上的 \citep{tang2023robust} 建立了内在联系：\BvM 定理证明了后验渐近正态，从而可以使用高斯目标分布的最优 MCMC 分析。
    \item \textbf{关键词:} MALA； 马尔可夫链蒙特卡洛； 混合时间； conductance； 高维推断
\end{enumerate}

\subsection{Minimax Optimal Rates for Regression on Manifolds and Distributions}
\begin{enumerate}
    \item \textbf{作者:} Rong Tang, Yun Yang
    \item \textbf{链接:} \url{https://arxiv.org/abs/2501.00001}
    \item \textbf{年份:} 2025
    \item \textbf{篇幅:} 待补充
    \item \textbf{核心贡献:}
        \begin{itemize}
            \item 在三类层次递进的设定下系统建立极小极大下界
            \item 提出 $\gamma$-Hölder IPM 统一框架，统一全变差距离（$\gamma=0$）和 Wasserstein 距离（$\gamma=1$）
            \item Regime 3 揭示 $d_X, d_Y, \beta_X, \beta_Y$ 等几何参数如何共同决定收敛速率
            \item 提出融合对抗学习与同步最小二乘的混合估计量
        \end{itemize}
    \item \textbf{摘要:} 本文将流形分布估计推广到条件设定，研究协变量依赖的流形上条件分布回归问题。通过建立 $\gamma$-Hölder IPM 统一框架，系统刻画了从经典密度回归到协变量依赖流形回归的完整极小极大理论。关键发现是：收敛速率仅依赖内在维度 $(d_X, d_Y)$ 而非环境维度 $(D_X, D_Y)$。
    \item \textbf{关键词:} 分布回归； 流形回归； 极小极大速率； IPM； 小波分解
\end{enumerate}

\subsection{Conditional Diffusion Models are Minimax-Optimal and Manifold-Adaptive}
\begin{enumerate}
    \item \textbf{作者:} Rong Tang, Lizhen Lin, Yun Yang
    \item \textbf{链接:} \url{https://arxiv.org/abs/2501.00002}
    \item \textbf年份:} 2025
    \item \textbf{篇幅:} 16 页
    \item \textbf{核心贡献:}
        \begin{itemize}
            \item 首次建立条件扩散模型的流形自适应理论
            \item 在经典密度回归设定下推广已有结果（覆盖 $\alpha_X\in[0,\infty)$ 而非仅 $[0,1]$）
            \item 在 Wasserstein 度量下，条件扩散模型的收敛速率仅依赖 $(d_X, d_Y)$
            \item 证明了条件扩散模型在一般条件设定下的极小极大最优性
        \end{itemize}
    \item \textbf{摘要:} 本文将 \citep{tang2024diffusion} 的流形自适应理论推广到条件生成设定。在两种场景下均达到极小极大最优收敛速率：经典密度回归（基于全变差度量）和高维回归（基于 Wasserstein 度量）。核心机制是条件分数函数的条件流形结构。
    \item \textbf{关键词:} 条件扩散模型； 流形自适应； 条件生成； 极小极大最优； Wasserstein 度量
\end{enumerate}

\section{补充参考资料}

\subsection{Nicolaescu: Introduction to Topology and Geometry}
\textbf{作者:} Liviu I. Nicolaescu
\textbf{年份:} 2024（第二版）
\textbf{篇幅:} 约 600 页

\textbf{内容简介:} 关于拓扑与几何的综合教材，涵盖流形理论、de Rham 上同调、黎曼几何等内容。是理解本笔记中黎曼几何概念的重要基础参考资料。

\textbf{用途:} 黎曼子流形、切空间、投影映射等几何概念的背景参考。

\subsection{Boumal: Introduction to Optimization on Smooth Manifolds}
\textbf{作者:} Nicolas Boumal
\textbf{年份:} 2023（第二版）
\textbf{篇幅:} 约 500 页

\textbf{内容简介:} 关于光滑流形上优化的权威教材。详细介绍了黎曼流形上的优化算法，包括 retraction、向量运输（vector transport）、信赖域方法等。是理解 RRWM 算法中 retraction 设计的重要参考。

\textbf{用途:} RRWM 算法的 retraction 和 proposal 设计参考（对应 \citep{tang2023robust}）。

\section{学习路径与前置基础知识}

根据论文的逻辑关系，建议学习顺序如下：

\begin{enumerate}
    \item \textbf{首先}: 学习 \citep{tang2022minimax}（Annals 2022）——
          理解流形分布估计的极小极大理论框架，建立理论基准
    \item \textbf{其次}: 学习 \citep{tang2024diffusion}（AISTATS 2024）——
          理解扩散模型如何达到极小极大最优，建立理论与实践的桥梁
    \item \textbf{然后}: 学习 \citep{tang2023robust}（JRSS B 2023）——
          理解流形上的贝叶斯推断与 BvM 定理，进入推断与计算主线
    \item \textbf{接着}: 学习 \citep{tang2024mala}（ICML 2024）——
          理解 MCMC 最优混合时间，将 BvM 与 MCMC 收敛分析连接
    \item \textbf{最后}: 学习 \citep{tang2025regression}（arXiv 2025）——
          理解条件设定下的极小极大理论，完成整条主线
\end{enumerate}

\textbf{前置基础知识}:
\begin{itemize}
    \item 黎曼几何基础：子流形、切空间、投影映射、retraction
    \item 经验过程理论：chaining、peeling、localization
    \item 极小极大理论：Fano 不等式、Varshamov-Gilbert 引理、Le Cam 方法
    \item 贝叶斯渐近理论：Bernstein-von Mises 定理
    \item 扩散模型基础：分数匹配、Langevin 动力学、Girsanov 定理
    \item MCMC 收敛分析：conductance、混合时间
\end{itemize}

\section{最终目标}

学习完该系列后，期望达到以下目标：

\begin{enumerate}
    \item 理解流形统计学习的完整理论体系：从极小极大基准到深度模型验证
    \item 掌握处理奇异分布（流形支撑）的核心技术：热核平滑、单位分解、Girsanov 定理
    \item 理解贝叶斯推断在流形约束下的渐近理论：\BvM 定理的几何版本
    \item 建立深度生成模型与极小极大理论之间的对应关系
    \item 能够独立开展流形统计学习方向的研究
\end{enumerate}
